\documentclass[10pt,a4paper]{amsart}
\usepackage{amsmath,amssymb,amsthm,graphicx}
% For commutative diagrams we use Xy-pic
\usepackage[all]{xy}

% We use 2cell for 2-commutative diagrams.
\xyoption{2cell}
% Theorem environments.
%
\theoremstyle{plain}
\newtheorem{theorem}[subsection]{Theorem}
\newtheorem{proposition}[subsection]{Proposition}
\newtheorem{lemma}[subsection]{Lemma}

\theoremstyle{definition}
\newtheorem{definition}[subsection]{Definition}
\newtheorem{example}[subsection]{Example}
\newtheorem{exercise}[subsection]{Exercise}
\newtheorem{situation}[subsection]{Situation}

\theoremstyle{remark}
\newtheorem{remark}[subsection]{Remark}
\newtheorem{remarks}[subsection]{Remarks}

\numberwithin{equation}{subsection}

% Macros
%
\def\lim{\mathop{\mathrm{lim}}\nolimits}
\def\colim{\mathop{\mathrm{colim}}\nolimits}
\def\Spec{\mathop{\mathrm{Spec}}}
\def\Hom{\mathop{\mathrm{Hom}}\nolimits}
\def\Ext{\mathop{\mathrm{Ext}}\nolimits}
\def\SheafHom{\mathop{\mathcal{H}\!\mathit{om}}\nolimits}
\def\SheafExt{\mathop{\mathcal{E}\!\mathit{xt}}\nolimits}
\def\Sch{\mathit{Sch}}
\def\Mor{\mathop{\mathrm{Mor}}\nolimits}
\def\Ob{\mathop{\mathrm{Ob}}\nolimits}
\def\Sh{\mathop{\mathit{Sh}}\nolimits}
\def\NL{\mathop{N\!L}\nolimits}
\def\CH{\mathop{\mathrm{CH}}\nolimits}
\def\proetale{{pro\text{-}\acute{e}tale}}
\def\etale{{\acute{e}tale}}
\def\QCoh{\mathit{QCoh}}
\def\Ker{\mathop{\mathrm{Ker}}}
\def\Im{\mathop{\mathrm{Im}}}
\def\Coker{\mathop{\mathrm{Coker}}}
\def\Coim{\mathop{\mathrm{Coim}}}

% Boxtimes
%
\DeclareMathSymbol{\boxtimes}{\mathbin}{AMSa}{"02}

%
% Macros for moduli stacks/spaces
%
\def\QCohstack{\mathcal{QC}\!\mathit{oh}}
\def\Cohstack{\mathcal{C}\!\mathit{oh}}
\def\Spacesstack{\mathcal{S}\!\mathit{paces}}
\def\Quotfunctor{\mathrm{Quot}}
\def\Hilbfunctor{\mathrm{Hilb}}
\def\Curvesstack{\mathcal{C}\!\mathit{urves}}
\def\Polarizedstack{\mathcal{P}\!\mathit{olarized}}
\def\Complexesstack{\mathcal{C}\!\mathit{omplexes}}
% \Pic is the operator that assigns to X its picard group, usage \Pic(X)
% \Picardstack_{X/B} denotes the Picard stack of X over B
% \Picardfunctor_{X/B} denotes the Picard functor of X over B
\def\Pic{\mathop{\mathrm{Pic}}\nolimits}
\def\Picardstack{\mathcal{P}\!\mathit{ic}}
\def\Picardfunctor{\mathrm{Pic}}
\def\Deformationcategory{\mathcal{D}\!\mathit{ef}}
\usepackage{cleveref}
\crefname{originalproposition}{original proposition}{original propositions}
\Crefname{originalproposition}{Original proposition}{Original propositions}
\makeatletter
\newlabel{proposition-etale-morphisms}{{03PC}{source}{}{sourcepointer-1}{}}
\newlabel{proposition-etale-morphisms@cref}{{[originalproposition][0][]03PC}{[1][0][]source}{}{}{}}
\makeatother
\title{A geometric point of an etale neighbourhood lifts to some member of any etale covering}
\author{The Stacks Project Authors}
\date{stacks-project:etale-cohomology.tex:local-snapshot-20261011; GNU FDL 1.2 or later}
\begin{document}
\maketitle
\noindent\textit{Editorial note.} All delivered mathematical text is the author's own, copied from the pinned Stacks Project chapter file etale-cohomology.tex; no statement or proof step was rewritten, reordered or removed.\par\medskip
\noindent\textit{Editorial note.} The preamble is the chapter's own native preamble excerpt used for this source: the Xy-pic setup at lines 14-18 and the theorem declarations with the mathematical macros at lines 154-217 of sources/stacks-etale-cohomology/preamble.tex. The Xy-pic xy environment with its curved arrows is required because the counted proof displays a cartesian square with the section sigma drawn as a curved arrow. The quoted class line and the xr-hyper externaldocument block are not included because the assembler supplies a plain amsart wrapper; no macro was invented or redefined.\par\medskip
\noindent\textit{Editorial note.} The same-chapter referenced proposition is located at its \textbackslash{}label line inside the pinned chapter file.\par\medskip
\noindent\textit{Allowed standard prerequisites (editorial).}\begin{itemize}
\item Etale neighbourhood of a geometric point, and etale covering of a scheme. Same chapter, the definition of an etale neighbourhood (U, ubar) with U etale over S, and the standard notion of an etale covering family. Supplies the objects and the surjectivity U = union of the images of the U\_i; the construction of the lift is the counted content.
\item Structure of etale morphisms over an algebraically closed field. Same chapter, Proposition labelled proposition-etale-morphisms (source line 2592, Stacks tag 03PC); invoked at proof lines 3329-3333. Supplies that a base change of an etale morphism is etale and that a scheme etale over the spectrum of an algebraically closed field is a disjoint union of trivial extensions; the section argument and the lifting conclusion are retained in the counted span.
\end{itemize}\medskip
\ifdefined\hypertarget\hypertarget{sourcepointer-1}{}\fi
\noindent\textit{Source reference.} Original proposition 03PC: sources/stacks-etale-cohomology/pinned-original-source.tex, source lines 2592--2592. Same-chapter pointer invoked twice, at proof lines 3330 and 3333, for base change of etale morphisms and for the shape of a scheme etale over an algebraically closed field; the cartesian square, the nonemptyness argument and the section construction are retained verbatim in the counted span.\par\medskip
\begin{lemma}
\label{lemma-geometric-lift-to-cover}
Let $S$ be a scheme.
Let $\overline{s}$ be a geometric point of $S$.
Let $(U, \overline{u})$ be an \'etale neighborhood of $\overline{s}$.
Let $\mathcal{U} = \{\varphi_i : U_i \to U \}_{i\in I}$ be an \'etale covering.
Then there exist $i \in I$ and $\overline{u}_i : \overline{s} \to U_i$
such that $\varphi_i : (U_i, \overline{u}_i) \to (U, \overline{u})$
is a morphism of \'etale neighborhoods.
\end{lemma}

\begin{proof}
As $U = \bigcup_{i\in I} \varphi_i(U_i)$, the fibre product
$\overline{s} \times_{\overline{u}, U, \varphi_i} U_i$ is not empty
for some $i$. Then look at the cartesian diagram
$$
\xymatrix{
\overline{s} \times_{\overline{u}, U, \varphi_i} U_i
\ar[d]^{\text{pr}_1} \ar[r]_-{\text{pr}_2} & U_i
\ar[d]^{\varphi_i} \\
\Spec(k) = \overline{s} \ar@/^1pc/[u]^\sigma
\ar[r]^-{\overline{u}} & U
}
$$
The projection $\text{pr}_1$ is the base change of an \'etale morphisms so it
is \'etale, see
Proposition \ref{proposition-etale-morphisms}.
Therefore, $\overline{s} \times_{\overline{u}, U, \varphi_i} U_i$
is a disjoint union of finite separable extensions of $k$, by
Proposition \ref{proposition-etale-morphisms}. Here
$\overline{s} = \Spec(k)$. But $k$ is algebraically closed, so all
these extensions are trivial, and there exists a section $\sigma$ of
$\text{pr}_1$. The composition
$\text{pr}_2 \circ \sigma$ gives a map compatible with $\overline{u}$.
\end{proof}

\end{document}
